\section{Version-by-version record}\label{sec:versions}

Table~\ref{tab:versions} lists every change from v0.3 to v0.9, the validation evidence and the decision. The
validation numbers are the score of \S\ref{sec:protocol}, where lower is better. Changes that failed the rule are
reported in the same way as changes that passed.

\begin{table*}[t]
\caption{Every change, the validation evidence, and the decision. Val: validation MAE in units of network spread
(lower is better); $\pm$: seed standard deviation of P0. Test effects are reported in \S\ref{sec:results}.}
\label{tab:versions}
\small
\begin{tabular}{@{}lp{5.9cm}p{6.3cm}l@{}}
\toprule
Version & Change & Validation evidence & Decision \\
\midrule
v0.3 & First release: three leagues, pure model with candidates P1 (wider), P2 (no graph), P3 (sparse experts);
hybrid H1 with IFS 0.25 and AIFS runs at least 48 h old; STRICT rule only &
WMO: removing the graph is worse (0.3948 vs 0.3864) & P0 kept \\
\addlinespace
v0.4 & ISSUED-BEFORE rule; four models in the hybrid; learned-correction hybrid H2; bias-corrected four-model
ensemble added as an opponent; cross-network training P4 (airports $\leftrightarrow$ WMO); fixed position scaling &
H2 vs H1: METAR 0.348 vs 0.372; WMO 0.3732 vs 0.3783; CPCB 0.4710 vs 0.4540.
P4: METAR 0.3920 vs P0 0.3934 $\pm$ 0.0012; WMO 0.3889 vs 0.3879 $\pm$ 0.0010.
CPCB: P1 0.4797 vs P0 0.4922 $\pm$ 0.0038 (v0.6 rerun) &
\begin{tabular}[t]{@{}l@{}}H2 (METAR, WMO)\\ H1 (CPCB)\\ P4 rejected\\ P1 (CPCB)\end{tabular} \\
\addlinespace
v0.5 & Hybrid H3: H2 plus cloud cover, solar radiation, dew point, pressure and wind direction from the same runs &
METAR 0.3539 vs H2 0.3481; CPCB (in the v0.6 run) 0.4932 vs H1 0.4586 & rejected \\
\addlinespace
v0.6 & Cross-network training for CPCB, with the 20 WMO stations as partner (hours before CPCB validation only);
fix to the foundation-model runner & CPCB: P4 0.4735 vs P0 0.4922 $\pm$ 0.0038 (P1 0.4797) & P4 (CPCB) \\
\addlinespace
v0.7 & Weekly rolling fine-tuning on hours observed before each week, simulated on validation &
CPCB: rolling 0.4727 vs frozen 0.4650 & rejected \\
\addlinespace
v0.8 & 14-day encoder (P6), weight averaging (P7), both (P8), all on top of P4 &
CPCB: P6 0.4714, P7 0.4721, P8 0.4678 vs P0 0.4922 $\pm$ 0.0038 & P8 (CPCB) \\
\addlinespace
v0.9 & 21-day patch encoder PWM-X (P9), plus masked-patch reconstruction (P10), on P4 with weight averaging;
data-threshold study (\S\ref{sec:v09}) &
CPCB: P9 0.4606, P10 0.4597 vs P8 0.4678, P0 0.4922 $\pm$ 0.0038 & P10 (CPCB) \\
\bottomrule
\end{tabular}
\end{table*}

\paragraph{What each version changed in practice.}
\begin{itemize}
\item \textbf{v0.3} set up the leagues. It compared global models only under STRICT, which our own analysis showed
to favour them, because their latest run often postdates our forecast start.
\item \textbf{v0.4} added the ISSUED-BEFORE rule and the learned-correction hybrid. At the airports this moved the
hybrid's MAE from 1.170 to 1.008 °C (temperature), 5.94 to 5.49\% (humidity) and 0.962 to 0.889 m/s (wind).
Cross-network training between the airports and the WMO stations did not pass the rule, and was not used there.
\item \textbf{v0.5} added more predictors from the global runs. It failed in both leagues where it was tried. A
fresh retraining of v0.4's design reproduced v0.4's test numbers exactly, which confirms that the pipeline is
deterministic.
\item \textbf{v0.6} targeted the CPCB league, the one league where the pure model lost to the foundation models.
That league has 16 months of training data. Cross-network training with the WMO stations improved validation by
about $5\sigma$. The same version fixed a bug that had invalidated v0.4's TimesFM 2.5 wind comparison on CPCB:
after running out of GPU memory on wind, the runner had reused the humidity forecast. We withdrew that number.
\item \textbf{v0.7} tested weekly retraining. Simulated on the validation period, it made the forecast worse, and it
was rejected.
\item \textbf{v0.8} let the model read the 14 days it already received hour by hour instead of only as an average,
and averaged its weights. On one seed, each change alone improved on P4 by less than $1\sigma$; together (P8)
they scored $6.4\sigma$ better than P0, most of which is P4's own gain. P8's
margin over v0.6's P4 (0.0057) is smaller than $2\sigma$, as was P4's margin over v0.4's P1 (0.0062). The rule
compares every candidate with P0, so its preference among passing candidates is not decisive.
\end{itemize}
Versions 0.6--0.9 were run on CPCB only (the threshold study also covers the airports). The airport and WMO results below use v0.4's chosen designs: P0 for the
pure model and H2 for the hybrid.
